\begin{proof}
Let $P_E,Q_E$ be the exact joint error Gramians. Subtracting the
reference Stein equations gives
$P_E-P_0=(I-\Phi_T)^{-1}D_P$, and the dual identity for $Q_E$.
The identity forcings have solutions
$\Udiag(X_A,X_B)$ and $\Udiag(Y_A,Y_B)$, respectively.
Lemma~\ref{I:U:lem:positive-resolvent} therefore gives
\[
 0\preceq P_E\preceq P_0+e_PI=:\overline P,\qquad
 0\preceq Q_E\preceq Q_0+e_QI=:\overline Q.
\]
For positive semidefinite matrices the quantity
$\lambda_{\max}(P^{1/2}QP^{1/2})$ is separately monotone in $P,Q$:
use congruence for $Q$, and the equality of the nonzero spectra of
$P^{1/2}QP^{1/2}$ and $Q^{1/2}PQ^{1/2}$ for $P$.
Apply Theorem~\ref{I:U:thm:exact}. The reference matrices commute.
On the retained symmetric and antisymmetric subspaces their pairs
of eigenvalues are $(2\sigma_j,0)$ and $(0,2\sigma_j)$; on a tail
coordinate they are $(\sigma_j,\sigma_j)$.
The largest eigenvalue of $\overline P\overline Q$ is exactly the
right-hand side of~\eqref{I:U:eq:reference-truncation}.

For the explicit estimates put
$S=\sum_aA_{a,12}\Sigma_2A_{a,12}^*$ and
$U=\sum_aA_{a,22}\Sigma_2A_{a,12}^*$.
Direct multiplication of the balanced Stein equations gives
\[
 D_P=-\begin{pmatrix}0&0&S\\0&0&U\\S&U^*&S\end{pmatrix}.
\]
Here $\|S\|\le M_P\delta$, and operator Cauchy--Schwarz, together
with $\sum_aA_{a,22}\Sigma_2A_{a,22}^*\preceq\Sigma_2$, gives
$\|U\|\le\sqrt{M_P}\delta$.
Split $D_P$ into its $S$ and $U$ parts; their norms are at most
$2\|S\|$ and $\|U\|$. Hence
$\|D_P\|\le(2M_P+\sqrt{M_P})\delta$.
Apply the same calculation to $A_a^*,c$, followed by changing the
sign of the last joint block, to obtain the asserted $D_Q$ bound.
\end{proof}
